% Generated draft: requires mathematical and source review.
\documentclass[11pt]{article}
\usepackage[margin=25mm]{geometry}
\usepackage{fontspec}
\setmainfont{Latin Modern Roman}
\usepackage{amsmath,amssymb,mathtools,unicode-math}
\setmathfont{Latin Modern Math}
\usepackage{enumitem,needspace,xurl,hyperref}
\setlength{\parindent}{0pt}
\setlength{\parskip}{5pt}
\setlength{\emergencystretch}{3em}
\newcommand{\N}{\mathbb{N}}
\newcommand{\Z}{\mathbb{Z}}
\newcommand{\Q}{\mathbb{Q}}
\newcommand{\R}{\mathbb{R}}
\newcommand{\C}{\mathbb{C}}
\newcommand{\F}{\mathbb{F}}
\newcommand{\A}{\mathbb{A}}
\newcommand{\PP}{\mathbb{P}}
\newcommand{\E}{\mathbb{E}}
\newcommand{\eps}{\varepsilon}
\newcommand{\dd}{\,\mathrm{d}}
\begin{document}
\section*{UnsolvedMath ID: 700627}
\textbf{Source problem code:} catalogue entry
\subsection{Definitions and mathematical statement}
Work in the first-order language of groups. A quasiidentity is a universally quantified sentence of the form
\[
\forall x_1\cdots\forall x_n\left(\bigwedge_{i=1}^{m}u_i(\mathbf{x})=v_i(\mathbf{x})\ \Longrightarrow\ u(\mathbf{x})=v(\mathbf{x})\right),
\]
where the expressions are group terms and the finite conjunction may be empty. For a group \(G\), let \(\operatorname{QId}(G)\) be the set of quasiidentities true in \(G\), and define its generated quasivariety by
\[
\operatorname{Qvar}(G)=\{H:H\text{ is a group and }H\models\operatorname{QId}(G)\}.
\]
Let \(\mathcal T\) denote the class of torsion-free groups: a group \(H\) belongs to \(\mathcal T\) if, for every \(h\in H\) and every integer \(r\geq1\), \(h^r=1\) implies \(h=1\). A group \(G\) is soluble if its derived series, defined by \(G^{(0)}=G\) and \(G^{(j+1)}=[G^{(j)},G^{(j)}]\), reaches the trivial group after finitely many steps. A free metabelian group on a set \(X\) is \(F(X)/F(X)''\), where \(F(X)\) is the free group on \(X\) and \(F(X)''=[[F(X),F(X)],[F(X),F(X)]]\).

The supplied question concerns every finitely generated torsion-free soluble group \(G\) such that, for some set \(X\) with \(|X|\geq2\),
\[
F(X)/F(X)''\in\operatorname{Qvar}(G).
\]
Here membership means that the quasivariety contains such a group as a member; no embedding in \(G\) is assumed. The source does not specify its rank beyond requiring non-abelianness.

\textbf{Question (A. I. Budkin, 1986).} For every such \(G\), does there exist an independent system \(\Sigma\) of quasiidentities defining \(\operatorname{Qvar}(G)\) within \(\mathcal T\), that is, satisfying
\[
\forall H\in\mathcal T,\qquad H\in\operatorname{Qvar}(G)\ \Longleftrightarrow\ H\models\Sigma\,?
\]
No finiteness requirement on \(\Sigma\) is stated.

The independence convention requires clarification. Independence relative to torsion-free groups means
\[
\forall\sigma\in\Sigma\ \exists H_\sigma\in\mathcal T:\qquad
H_\sigma\models\Sigma\setminus\{\sigma\},\qquad H_\sigma\not\models\sigma.
\]
Independence over all groups has the same condition with \(H_\sigma\) allowed to be any group. These are distinct requirements in general. The supplied wording specifies the class in which the system defines the quasivariety, but does not explicitly specify the class relative to which independence is intended. The displayed relative condition therefore records one precise reading, rather than a verified resolution of that source qualification.
\subsection{Short English statement}
Budkin asks whether every quasivariety generated by a finitely generated torsion-free soluble group, and containing a non-abelian free metabelian group, has an independent quasiidentity axiomatization within torsion-free groups. The supplied wording leaves unclear whether independence must hold relative to torsion-free groups or relative to all groups.
\subsection{Sources}
\begin{itemize}
\item Local UnsolvedMath registry, permanent ID \(700627\), exact source code \texttt{catalogue entry}, title \emph{Kourovka Notebook Problem 10.10}; supplied complete source record, upstream revision \texttt{d371ee917ac9647fbb06cd6fc529deec6f3dda82}.
\item Supplied attribution: A. I. Budkin, proposed in 1986; original problem number 10.10 in \emph{The Kourovka Notebook (Unsolved Problems in Group Theory)}, 10th issue (1986). The record also identifies the 21st issue edition (2026), \texttt{arXiv:1401.0300v48}. Its collection description names E. I. Khukhro and V. D. Mazurov as editors.
\item The registry labels the problem \texttt{open} and records an update timestamp of \texttt{2026-10-16T14:00:00Z}. This is supplied registry information, not an independently verified assessment of the literature or present mathematical status. The cited Notebook edition has not been independently consulted for this draft.
\end{itemize}
\end{document}
